% Appendix B: Theoretical Analysis of Cross-Attention Design
% For MoM-Driver NeurIPS 2026 Paper (Compressed Version)

\section{Theoretical Analysis of Cross-Attention Design}
\label{appendix:cross_attention}

This appendix analyzes our State-Frozen cross-attention mechanism compared to Linear Interpolation Cross-Attention (LICA) used in prior work~\cite{lady}. \textbf{This analysis is not claimed as a main contribution}, but demonstrates our theoretical understanding of cross-attention design choices.

\subsection{Problem Setup}

In multi-modal linear attention, cross-attention between queries $\mathbf{Q} \in \mathbb{R}^{N_q \times d}$ and key-values $\{\mathbf{K}, \mathbf{V}\} \in \mathbb{R}^{N_{kv} \times d}$ is typically implemented via \textbf{concatenation}:

\begin{enumerate}
    \item \textbf{LICA approach}: $\mathbf{X} = [\mathbf{K}; \mathbf{V}; \mathbf{Q}] \rightarrow \text{LinearAttn}(\mathbf{X}) \rightarrow \text{extract queries}$
    \item \textbf{State-Frozen approach}: Prevent query tokens from updating recurrent state
\end{enumerate}

For RWKV-style recurrence: $\mathbf{S}_t = \lambda \mathbf{S}_{t-1} + \mathbf{k}_t \mathbf{v}_t^\top$, $\mathbf{o}_t = \mathbf{q}_t \mathbf{S}_t$

\subsection{Systematic Biases in LICA}

\subsubsection{Bias 1: Positional Decay ($\varepsilon_{decay}$)}

\begin{theorem}[Positional Decay Bias]
In LICA, visual information at position $i$ is attenuated by $\lambda^{N_q}$ when accessed by query tokens, where $N_q$ is the number of queries.
\end{theorem}

\begin{proof}[Proof Sketch]
Visual token $(\mathbf{k}_i, \mathbf{v}_i)$ processed at position $i < N_{kv}$ contributes to query $\mathbf{q}_j$ at position $j = N_{kv} + \delta$:
\begin{align}
\mathbf{o}_j^{(i)} &= \mathbf{q}_j (\lambda^{\delta} \mathbf{k}_i \mathbf{v}_i^\top) = \lambda^{\delta} (\mathbf{q}_j \mathbf{k}_i) \mathbf{v}_i
\end{align}

The decay factor $\lambda^{\delta}$ artificially suppresses visual information. For typical values $\lambda = 0.9$, $N_q = 64$:
$$\varepsilon_{decay} = 1 - \lambda^{64} \approx 0.9998$$

This means visual information is almost entirely suppressed!
\end{proof}

\subsubsection{Bias 2: Query Contamination ($\varepsilon_{contam}$)}

\begin{theorem}[Query Contamination Bias]
In LICA, later queries read state contaminated by earlier queries, violating cross-attention semantics.
\end{theorem}

\begin{proof}[Proof Sketch]
Query $\mathbf{q}_j$ updates state: $\mathbf{S}_j = \lambda \mathbf{S}_{j-1} + \mathbf{k}_j \mathbf{v}_j^\top$

Later query $\mathbf{q}_{j+1}$ reads contaminated state:
$$\mathbf{o}_{j+1} = \mathbf{q}_{j+1} \mathbf{S}_{j+1} = \mathbf{q}_{j+1}(\lambda \mathbf{S}_j + \mathbf{k}_{j+1} \mathbf{v}_{j+1}^\top)$$

The contamination term $\lambda (\mathbf{q}_{j+1} \mathbf{k}_j) \mathbf{v}_j$ means queries influence each other's outputs.
\end{proof}

\subsection{State-Frozen Solution}

\begin{definition}[State-Frozen Cross-Attention]
For query positions $t \geq N_{kv}$:
\begin{align}
\text{State:} \quad &\mathbf{S}_t = \mathbf{S}_{t-1} \quad \text{(frozen, no update)} \\
\text{Output:} \quad &\mathbf{o}_t = \mathbf{q}_t \mathbf{S}_t \quad \text{(normal read)}
\end{align}
\end{definition}

\begin{theorem}[Bias Elimination]
State-Frozen eliminates both biases: $\varepsilon_{decay} = 0$ and $\varepsilon_{contam} = 0$.
\end{theorem}

\begin{proof}
\textbf{No decay}: State remains $\mathbf{S}_{N_{kv}}$ for all queries, eliminating positional decay.

\textbf{No contamination}: Each query reads identical pure visual state $\mathbf{S}_{N_{kv}}$:
$$\mathbf{o}_j = \mathbf{q}_j \mathbf{S}_{N_{kv}} = \mathbf{q}_j \left(\sum_{i=1}^{N_{kv}} \lambda^{N_{kv}-i} \mathbf{k}_i \mathbf{v}_i^\top \right)$$
\end{proof}

\subsection{Empirical Observations}

Despite theoretical biases, LADY achieves PDMS = 90.9 with LICA, suggesting:
\begin{itemize}
    \item \textbf{Network compensation}: Deep networks may adapt parameters to counteract biases
    \item \textbf{Task robustness}: Some driving scenarios may tolerate attention artifacts
    \item \textbf{Implicit regularization}: Biases might act as beneficial constraints
\end{itemize}

\subsection{Implementation}

In our MoM blocks, State-Frozen is controlled via \texttt{frozen\_from} parameter:

\begin{verbatim}
# LICA baseline
frozen_from = -1  # No freezing, standard behavior

# State-Frozen (our approach)  
frozen_from = kv_len  # Freeze state updates from query positions
\end{verbatim}

\textbf{Computational cost}: Identical to LICA ($O(N_{total} \cdot d^2)$), only state update logic differs.

\subsection{Connection to Standard Cross-Attention}

\begin{theorem}[Transformer Equivalence]
With $\lambda = 1$ and infinite precision, State-Frozen recovers standard Transformer cross-attention (without softmax).
\end{theorem}

\begin{proof}
$$\mathbf{o}_j = \mathbf{q}_j \left(\sum_{i=1}^{N_{kv}} \mathbf{k}_i \mathbf{v}_i^\top \right) = \sum_{i=1}^{N_{kv}} (\mathbf{q}_j^\top \mathbf{k}_i) \mathbf{v}_i$$

This matches linear attention approximation of Transformer cross-attention.
\end{proof}

\subsection{Key Insights}

\begin{itemize}
    \item \textbf{LICA introduces systematic biases} that can be quantified ($\varepsilon_{decay} \approx 0.9998$ for typical parameters)
    \item \textbf{State-Frozen provides principled bias elimination} while maintaining efficiency
    \item \textbf{Both approaches work in practice}, demonstrating robustness of deep learning
    \item \textbf{Design choice depends on priorities}: theoretical soundness vs. empirical optimization
\end{itemize}

While we implement State-Frozen for theoretical rigor, we acknowledge both formulations can be effective depending on task requirements and network capacity for bias compensation.